%%%%%%%%%%%%%%%%%
% This is an sample CV template created using altacv.cls
% (v1.3, 10 May 2020) written by LianTze Lim (liantze@gmail.com). Now compiles with pdfLaTeX, XeLaTeX and LuaLaTeX.
% (v1.7.1b, 11 Jan 2024) forked by Nicolás Omar González Passerino (nicolas.passerino@gmail.com)
%
%% It may be distributed and/or modified under the
%% conditions of the LaTeX Project Public License, either version 1.3
%% of this license or (at your option) any later version.
%% The latest version of this license is in
%%    http://www.latex-project.org/lppl.txt
%% and version 1.3 or later is part of all distributions of LaTeX
%% version 2003/12/01 or later.
%%%%%%%%%%%%%%%%

%% If you need to pass whatever options to xcolor
\PassOptionsToPackage{dvipsnames}{xcolor}

%% If you are using \orcid or academicons
%% icons, make sure you have the academicons
%% option here, and compile with XeLaTeX
%% or LuaLaTeX.
% \documentclass[10pt,a4paper,academicons]{altacv}

%% Use the "normalphoto" option if you want a normal photo instead of cropped to a circle
% \documentclass[10pt,a4paper,normalphoto]{altacv}

%% Fork (before v1.6.5a): CV dark mode toggle enabler to use a inverted color palette.
%% Use the "darkmode" option if you want a color palette used to 
% \documentclass[10pt,a4paper,ragged2e,withhyper,darkmode]{altacv}

\documentclass[10pt,a4paper,ragged2e,withhyper]{altacv}

%% AltaCV uses the fontawesome5 and academicons fonts
%% and packages.
%% See http://texdoc.net/pkg/fontawesome5 and http://texdoc.net/pkg/academicons for full list of symbols. You MUST compile with XeLaTeX or LuaLaTeX if you want to use academicons.

%% Fork v1.6.5c: Overwriting sloppy environment to ignore any spaces and be used to solve overlapping cvtags
\newenvironment{sloppypar*}{\sloppy\ignorespaces}{\par}

% Change the page layout if you need to
\geometry{left=1.2cm,right=1.2cm,top=1cm,bottom=1cm,columnsep=0.75cm}

% The paracol package lets you typeset columns of text in parallel
\usepackage{paracol}

% Change the font if you want to, depending on whether
% you're using pdflatex or xelatex/lualatex
\ifxetexorluatex
  % If using xelatex or lualatex:
  \setmainfont{Roboto Slab}
  \setsansfont{Lato}
  \renewcommand{\familydefault}{\sfdefault}
\else
  % If using pdflatex:
  \usepackage[rm]{roboto}
  \usepackage[defaultsans]{lato}
  % \usepackage{sourcesanspro}
  \renewcommand{\familydefault}{\sfdefault}
\fi

% Fork (before v1.6.5a): Change the color codes to test your personal variant on any mode
\ifdarkmode%
    \definecolor{PrimaryColor}{HTML}{0A2540}   % Deep aerospace blue
    \definecolor{SecondaryColor}{HTML}{1F4E79} % Lighter blue accents
    \definecolor{ThirdColor}{HTML}{FF7A00}     % Accent orange (small highlights)
    \definecolor{BodyColor}{HTML}{333333}      % Main text
    \definecolor{EmphasisColor}{HTML}{111111}  % Strong text
    \definecolor{BackgroundColor}{HTML}{FFFFFF} % White background
\else%
  \definecolor{PrimaryColor}{HTML}{001F5A} % other blue shit important
  \definecolor{SecondaryColor}{HTML}{0039AC} % blue shit important
  \definecolor{ThirdColor}{HTML}{000000} % lines
  \definecolor{BodyColor}{HTML}{666666} % body text
  \definecolor{EmphasisColor}{HTML}{000000} % titles
  \definecolor{BackgroundColor}{HTML}{FFFFFF}
\fi%

\colorlet{name}{PrimaryColor}
\colorlet{tagline}{SecondaryColor}
\colorlet{heading}{PrimaryColor}
\colorlet{headingrule}{ThirdColor}
\colorlet{subheading}{SecondaryColor}
\colorlet{accent}{SecondaryColor}
\colorlet{emphasis}{EmphasisColor}
\colorlet{body}{BodyColor}
\pagecolor{BackgroundColor}

% Change some fonts, if necessary
\renewcommand{\namefont}{\Huge\rmfamily\bfseries}
\renewcommand{\personalinfofont}{\small\bfseries}
\renewcommand{\cvsectionfont}{\Large\bfseries}
\renewcommand{\cvsubsectionfont}{\large\bfseries}

% Change the bullets for itemize and rating marker
% for \cvskill if you want to
\renewcommand{\itemmarker}{{\small\textbullet}}
\renewcommand{\ratingmarker}{\faCircle}

%% sample.bib contains your publications
%% \addbibresource{main.bib}

\begin{document}
    \name{Martim Silva}
    \tagline{Software, Telecommunications \& Computer Engineering Student | Team Leader in CubeSat, Robotics, and Software Projects }
    
    %% You can add multiple photos on the left or right
    \photoL{4cm}{AI_good_face}

    \personalinfo{
        \email{Martim.Alexandre.Dras@gmail.com}\smallskip
        \phone{+351-936-495-912}
        \location{Lisbon, Portugal}\\
        \linkedin{martim-silva-tech-dev/}
        \github{XanderDras}
        
        %\homepage{nicolasomar.me}
        %\medium{nicolasomar}
        %% You MUST add the academicons option to \documentclass, then compile with LuaLaTeX or XeLaTeX, if you want to use \orcid or other academicons commands.
        % \orcid{0000-0000-0000-0000}
        %% You can add your own arbtrary detail with
        %% \printinfo{symbol}{detail}[optional hyperlink prefix]
        % \printinfo{\faPaw}{Hey ho!}[https://example.com/]
        %% Or you can declare your own field with
        %% \NewInfoFiled{fieldname}{symbol}[optional hyperlink prefix] and use it:
        % \NewInfoField{gitlab}{\faGitlab}[https://gitlab.com/]
        % \gitlab{your_id}
    }
    
    \makecvheader
    %% Depending on your tastes, you may want to make fonts of itemize environments slightly smaller
    % \AtBeginEnvironment{itemize}{\small}
    
    %% Set the left/right column width ratio to 6:4.
    \columnratio{0.22}

    % Start a 2-column paracol. Both the left and right columns will automatically
    % break across pages if things get too long.
    \begin{paracol}{2}
        % ----- TECH STACK -----
        \cvsection{SKILL STACK}
            %% Fork v1.6.5c: The sloppypar* environment is used to avoid tags overlapping with section width
            \begin{sloppypar*}
                %% Fork 1.7.1b: Now in case you want to highlight any tag, just add a '/true' property next to its text and it will change tag's text and border colors.
                \cvtags{Team Leadership/true, CubeSat Design/true, Robotics/true, Embedded Systems/true, Python/true, C++/true, JavaScript/true, Node.js, Raspberry Pi, OpenCV, OOP, Git, Linux}
                \medskip

            \end{sloppypar*}
        % ----- TECH STACK -----
        
        % ----- LEARNING -----
        %\cvsection{Learning}
         %   \begin{sloppypar*}
         %       \cvtags{Uno, Dos/true, Tres, Cuatro/true, Cinco, Seis/true, Siete, Ocho/true, Nueve, Diez/true}
         %       \medskip

         %       \cvtags{Rojo/true, Amarillo, Azul/true, Verde, Violeta/true, Naranja, Marron/true, Blanco, Gris/true, Negro}
        %    \end{sloppypar*}
        % ----- LEARNING -----
        
        % ----- LANGUAGES -----
        \cvsection{Languages}
            \cvlang{English}{C2}\\
            \divider
            
            \cvlang{Portuguese}{Native}
            \bigskip
            %% Yeah I didn't spend too much time making all the
            %% spacing consistent... sorry. Use \smallskip, \medskip,
            %% \bigskip, \vpsace etc to make ajustments.
        % ----- LANGUAGES -----

       % ----- EDUCATION -----
        \cvsection{Education}
        
        \textbf{BSc Software, Telecommunications \& Computer Engineering}\\
        \smallskip
        Instituto Superior de Engenharia de Lisboa (ISEL)\\
        \smallskip
        Sep 2023 -- Jul 2026
        
        \divider
        
        \textbf{Vocational Course in Information Systems Management and Programming}\\
        \smallskip
        Escola Secundária Gago Coutinho\\
        \smallskip
        2020 -- 2023\\
        Student Excellence Award – Municipality of Vila Franca de Xira
        % ----- EDUCATION -----
            
        % ----- REFERENCES -----
        \cvsection{References}
            \cvref{Prof.\ Rui Duarte}{ISEL}{}
            \divider

            \cvref{Boss\ Paulo Gouveia}{YEPKIT}{}
        % ----- REFERENCES -----
        
        
        % ----- MOST PROUD -----
        % \cvsection{Most Proud of}
        
        % \cvachievement{\faTrophy}{Fantastic Achievement}{and some details about it}\\
        % \divider
        % \cvachievement{\faHeartbeat}{Another achievement}{more details about it of course}\\
        % \divider
        % \cvachievement{\faHeartbeat}{Another achievement}{more details about it of course}
        % ----- MOST PROUD -----
        
        % \cvsection{A Day of My Life}
        
        % Adapted from @Jake's answer from http://tex.stackexchange.com/a/82729/226
        % \wheelchart{outer radius}{inner radius}{
        % comma-separated list of value/text width/color/detail}
        % \wheelchart{1.5cm}{0.5cm}{%
        %   6/8em/accent!30/{Sleep,\\beautiful sleep},
        %   3/8em/accent!40/Hopeful novelist by night,
        %   8/8em/accent!60/Daytime job,
        %   2/10em/accent/Sports and relaxation,
        %   5/6em/accent!20/Spending time with family
        % }
        
        % use ONLY \newpage if you want to force a page break for
        % ONLY the current column
        \newpage
        
        %% Switch to the right column. This will now automatically move to the second
        %% page if the content is too long.
        \switchcolumn
        
        % ----- ABOUT ME -----
        \cvsection{About Me}
            \begin{quote}
            Award-winning engineering student with hands-on experience in CubeSat and robotics competitions, leading multidisciplinary teams to deliver high-impact projects. Skilled in software development with measurable results in international and national competitions.
            \end{quote}
        % ----- ABOUT ME -----

        % ----- PROJECTS -----
        \cvsection{Engineering Projects}
            \cvevent{PocketQube – Modular COTS Satellite Platform}{Academic Engineering Project}{Mar 2026 -- Jul 2026 (ONGOING)}{ISEL}
            \begin{itemize}
            \item Low-cost modular Pocket satellite platform using commercial off-the-shelf components.
            \item Developing satellite subsystems including power management, ESP32-C3, and RF communications (433 MHz).
            \item Implementing embedded firmware with FreeRTOS for telemetry acquisition and transmission.
            \end{itemize}
            
            \cvevent{Smart Car – Web-Controlled IoT Project}{\cvreference}{Jan 2023 -- Jul 2023}{}
            \begin{itemize}
                \item Developed web-controllable smart car using Raspberry Pi and WebSockets, providing live video feed via OpenCV.
                \item Designed a custom API for real-time telemetry and control.
            \end{itemize}
            
            %\cvevent{Project 2 }{\cvreference{\faGitlab}{https://gitlab.com/user/repo}\cvreference{| \faGlobe}{https://project-demo.com/}}{Mm YYYY -- Mm YYYY}{}
            %\begin{itemize}
            %    \item Item 1
            %    \item Item 2
            %\end{itemize}
        % ----- PROJECTS -----
        
        % ----- EXPERIENCE -----
        \cvsection{Experience}
            \cvevent{Software Developer}{YEPKIT}{May 2023 -- Jul 2023}{Portugal (Hybrid)}
            \begin{itemize}
                \item Developed backend services using Node.js and Docker.
                \item Contributed to full-stack development and API design, enhancing system reliability.
                \item Software engineering best practices including version control (Git) and modular design.
            \end{itemize}
            \divider
            
            \cvevent{Team Leader | CubeSat Portugal Bootcamp}{ISEL-Space Club}{April 2025}{Lisbon, Portugal}
            \begin{itemize}
                \item Led 30-member team in CubeSat design and integration, to deliver a prototype.
                \item Coordinated embedded systems, communications, and mission planning with guidance from Portuguese Space Agency experts.
                \item Strengthened team technical skills and collaboration across hardware and software.
            \end{itemize}

            \cvevent{Team Leader | Robotics Competitions}{Escola Secundária Gago Coutinho}{2022 -- 2023}{Portugal}
            \begin{itemize}
                \item Led teams to compete in International First Global Competition 2022 (Geneva, Switzerland), European RoboCup 2022, Festival Nacional da Robótica 2023, and CanSat 2022.
                \item Designed and programmed autonomous robots using Arduino, Python, and OpenCV.
                \item Oversaw project documentation, software versioning (Git), and cross-functional collaboration.
            \end{itemize}
        % ----- EXPERIENCE -----

        
        
        % ----- EDUCATION -----
        %\cvsection{Education}
         %   \cvevent{BSc Computer Systems Networking and Telecommunications}{Instituto Superior de Engenharia de Lisboa (ISEL)}{Sep 2023 -- Jul 2026}{Lisbon, Portugal}
         %   \begin{itemize}
          %     \item Skills developed: GitHub, Assembly, Python, Embedded Systems, Robotics
            %\end{itemize}
           % \divider

            
           %\cvevent{Associate Degree in Computer Programming, Specific Applications}{Escola Secundária Gago Coutinho}{2020 -- Jul 2023}{Portugal}
            %\begin{itemize}
             %   \item Team Leader in International and national robotics competitions and CanSat 2022
              %  \item Awarded Student Excellence by Municipal of Vila Franca de Xira
            %\end{itemize}
        % ----- EDUCATION -----
        
        
    \end{paracol}
\end{document}